
\documentclass[11pt]{book}
% GLAVNI FAJL. Tekst se piše u sections/*.tex; ovaj fajl ih samo uključuje redom.
% Kompajlira se uvijek ovaj fajl (sections imaju „% !TEX root = ../OzicaneNeuronskeMreze.tex“, pa
% editor sam kompajlira glavni fajl i kad se snima sekcija).
% Kompajlirati XeLaTeX-om (ne pdfLaTeX-om). latexmk pokrenut iz ovog foldera čita
% .latexmkrc i sve pomoćne fajlove stavlja u build/.
% Fontovi Source Serif/Sans/Code Pro moraju biti instalirani; ako nisu, zakomentarisati
% linije 7, 8 i 9 u langencyc.sty.
% Putanje do slika se pišu u odnosu na ovaj folder (i iz sekcija), npr. slike/neuron.png.
% Pregled svih naredbi predloška (listinzi, slike, tabele, gramatike) je na dnu fajla,
% iza \end{document}. Original predloška: University/courses/programski-jezici-i-prevodioci/
% materials/assignments/seminarski-uzorak/template.tex
\usepackage{langencyc}
\usepackage{url} % linkovi u bibliografiji
\begin{document}


\setcounter{chapter}{4}
\setcounter{page}{1}


\lstset{language=VHDL}

\chapter{Ožičeno programiranje neuralnih mreža}
\begin{center}
    Aid Mustafić
\end{center}

% !TEX root = ../OzicaneNeuronskeMreze.tex
% UVOD (pasus bez naslova sekcije, odmah ispod naslova poglavlja i imena autora):
% šta je ožičeno programiranje neuralnih mreža.
% Nadovezuje se na uvod u dio „Ožičeno programiranje“  i proširuje ga.
% Ne ponavljati podsekciju „Neuromorfno računarstvo“ iz poglavlja Analogni računari,
% nego se na nju pozvati.

% !TEX root = ../OzicaneNeuronskeMreze.tex
\section{Formalni neuron i principi neuralnih mreža}
% Profesor nije naveo izvore za ovu sekciju: dodati vlastite
% (npr. McCulloch i Pitts 1943, Rosenblatt 1958, udžbenik o neuralnim mrežama).
% Kratko: ova sekcija samo uvodi pojmove koje ostatak poglavlja koristi.

\subsection{Model formalnog neurona}
% Ulazi, težine, težinska suma, prag, izlaz. Slika neurona nacrtana u TikZ-u.

\subsection{Aktivacijske funkcije}
% Prag (step), sigmoid, ReLU; zašto je bitno koju funkciju hardver može realizovati.

\subsection{Perceptron i pravilo učenja}
% Pravilo korekcije težina; učenje = mijenjanje težina, tj. „programa“ mreže.

\subsection{Višeslojne mreže}
% Slojevi, propagacija unaprijed, kratko o učenju propagacijom greške unazad.

\subsection{Operacija MAC kao osnova hardverske realizacije}
% Množenje i akumulacija (multiply-accumulate): operacija koju ubrzava sav hardver
% iz narednih sekcija. Ovdje je uvesti, kasnije se samo pozivati na nju.

% !TEX root = ../OzicaneNeuronskeMreze.tex
\section{Historijski počeci: analogni i elektromehanički hardver}
% Prve realizacije su koristile fizičke otpornike, motore i analogna kola, jer
% tadašnji računari nisu imali dovoljnu procesorsku snagu.

\subsection{Mark I Perceptron}
% 1957/1958, Frank Rosenblatt: prva fizička implementacija neuronske mreže.
% Težine = potenciometri koje okreću elektromotori.
% Izvori: wikiann, galileo.

\subsection{ADALINE i MADALINE}
% 1959/1962, Bernard Widrow i Ted Hoff; ADALINE = Adaptive Linear Neuron.
% Memistori = otpornici s memorijom.
% Provjeriti u izvorima: memistor se obično opisuje kao elektrohemijski element,
% a ADALINE se najčešće datira u 1960. Izvor: stanford.

% !TEX root = ../OzicaneNeuronskeMreze.tex
\section{Tranzicija ka VLSI i FPGA}
% 1980-e i 1990-e: razvoj VLSI tehnologije (Very Large Scale Integration), fokus se
% prebacuje na integrisana kola specifične namjene (ASIC) i programabilne logičke
% nizove (FPGA).

\subsection{Od von Neumannove arhitekture ka namjenskim VLSI sistemima}
% Izvor: james2017.

\subsection{Digitalne, analogne i hibridne realizacije}
% Izvor: schuman2017.

\subsection{Realizacija neurona na FPGA}
% Kratki VHDL listing neurona (MAC + aktivacija) u shortlisting okruženju.
% Izvori: ijecce, awesomehw.

\subsection{Optimizacija hardvera i potrošnje}
% Od stotina vati do milivata. Izvor: hwimpl.

% !TEX root = ../OzicaneNeuronskeMreze.tex
\section{Specijalizovana integrisana kola (AI ASIC i TPU)}
% Era dubokog učenja: potreba za masovnim matričnim množenjem (MAC operacije),
% pa nastaju namjenski akceleratori za duboko učenje.

\subsection{Matrično množenje i MAC operacije}
% Veza sa sekcijom 1 (operacija MAC); zašto je opšti procesor tu spor.

\subsection{Preslikavanje modela na crossbar nizove}
% Kako se matematički model prebacuje na hardverske matrice otpornika i crossbar
% nizove. Izvor: gopalakrishnan2019 (koristi i za sekciju 5).

\subsection{Neuročipovi: IBM ZISC i CNAPS}
% Izvor: ijecce.

\subsection{Google TPU}
% Izvor: ijecce.

\subsection{Hardverska ograničenja}
% Preciznost, memorija, propusnost, potrošnja; šta arhitekture rade da ih zaobiđu.

% !TEX root = ../OzicaneNeuronskeMreze.tex
\section{Neuromorfni čipovi}
% Za razliku od klasičnih akceleratora, imitiraju biološki rad mozga u realnom
% vremenu: obrada vođena događajima (event-driven) i impulsne mreže (SNN).

\subsection{Obrada vođena događajima i impulsne neuronske mreže}
% Impulsni neuron, zašto event-driven obrada štedi energiju.

\subsection{Primjeri čipova: IBM TrueNorth, Intel Loihi, SpiNNaker}
% Izvor: openneuro.

\subsection{Otvoreni hardverski projekti}
% Projekti s kodom u VHDL/Verilog/SystemVerilog, digitalni i asinhroni procesori.
% Izvor: awesomehw.

\subsection{Memristori i RRAM}
% Buduće tehnologije: sinaptičke težine emulirane na atomskom nivou. Izvor: river.

% OPCIONALNO
%% !TEX root = ../OzicaneNeuronskeMreze.tex
% OPCIONO, samo ako prof. Ribić odobri na konsultacijama (po uzoru na poglavlje
% Analogni računari, sekcije „Veza sa programskim jezicima“).
% Uključuje se odkomentarisanjem \input{sekcije/06-veza-s-programskim-jezicima}
% u OzicaneNeuronskeMreze.tex.
\section{Veza s programskim jezicima}

\subsection{Šta je program kod neuralnog hardvera}
% Položaji potenciometara → otpor memistora → HDL i bitstream → prevedeni graf na
% TPU → impulsna mreža preslikana na neuromorfni čip.

\subsection{Prevođenje: od modela do sklopa}


% !TEX root = ../OzicaneNeuronskeMreze.tex
% Bibliografija: inline citati nisu obavezni, ali navesti sve izvore.
% Prijedlozi za dopunu (nisu od profesora): McCulloch i Pitts 1943; Rosenblatt 1958;
% Widrow i Hoff 1960 (Adaptive switching circuits); Widrow i Lehr 1990 (30 years of
% adaptive neural networks); Mead 1990 (Neuromorphic electronic systems);
% Misra i Saha 2010; Sze i dr. 2017; Jouppi i dr. 2017 (TPU); Merolla i dr. 2014
% (TrueNorth); Davies i dr. 2018 (Loihi); Furber i dr. 2014 (SpiNNaker).
% Linkovi za većinu: ../references/reading-list.md
\begin{thebibliography}{11}

    \bibitem{wikiann}
    Wikipedia contributors,
    \textit{History of artificial neural networks},
    Wikipedia -- The Free Encyclopedia.
    Dostupno na: \url{https://en.wikipedia.org/wiki/History_of_artificial_neural_networks}
    (pristupljeno: novembar 2026.)

    \bibitem{galileo}
    Galileo Unbound (blog),
    \textit{A short history of neural networks}, 2025.
    Dostupno na: \url{https://galileo-unbound.blog/2025/02/05/a-short-history-of-neural-networks/}

    \bibitem{stanford}
    Stanford University, CS Sophomore College,
    \textit{Neural Networks -- History}.
    Dostupno na: \url{https://cs.stanford.edu/people/eroberts/courses/soco/projects/neural-networks/History/history1.html}

    \bibitem{james2017}
    C. D. James i dr.,
    \textit{A historical survey of algorithms and hardware architectures for neural-inspired
        and neuromorphic computing applications},
    Biologically Inspired Cognitive Architectures, 2017.
    Dostupno na: \url{https://www.sciencedirect.com/science/article/am/pii/S2212683X16300561}

    \bibitem{schuman2017}
    C. D. Schuman i dr.,
    \textit{A Survey of Neuromorphic Computing and Neural Networks in Hardware},
    arXiv:1705.06963, 2017.
    Dostupno na: \url{https://arxiv.org/abs/1705.06963}

    \bibitem{hwimpl}
    % TODO: autori i godina (vidjeti na ResearchGate stranici)
    \textit{Hardware Implementation of Neural Networks}, ResearchGate.
    Dostupno na: \url{https://www.researchgate.net/publication/385988727_Hardware_Implementation_of_Neural_Networks}

    \bibitem{gopalakrishnan2019}
    R. Gopalakrishnan, Y. Chua, A. J. S. Kumar,
    \textit{Hardware-friendly Neural Network Architecture for Neuromorphic Computing},
    arXiv:1906.08853, 2019.
    Dostupno na: \url{https://arxiv.org/abs/1906.08853}

    \bibitem{ijecce}
    % TODO: autori, broj časopisa i godina (vidjeti u PDF-u)
    \textit{Hardware Implementation of Neural Network Using VHDL},
    International Journal of Electronics Communication and Computer Engineering (IJECCE).
    Dostupno na: \url{https://ijecce.org/administrator/components/com_jresearch/files/publications/IJECCE_671_1_FInal.pdf}

    \bibitem{openneuro}
    Open Neuromorphic,
    \textit{Neuromorphic Hardware Guide}.
    Dostupno na: \url{https://open-neuromorphic.org/neuromorphic-computing/hardware/}

    \bibitem{awesomehw}
    Open Neuromorphic,
    \textit{Awesome Neuromorphic HW}, GitHub repozitorij.
    Dostupno na: \url{https://github.com/open-neuromorphic/awesome-neuromorphic-hw}

    \bibitem{river}
    % TODO: autori i naslov knjige (vidjeti u PDF-u poglavlja)
    \textit{Technology and Hardware for Neuromorphic Computing}, poglavlje u knjizi,
    River Publishers.
    Dostupno na: \url{https://www.riverpublishers.com/downloadchapter.php?file=RP_9788770226103C2.pdf}

\end{thebibliography}


\end{document}

%%%%%%%%%%%%%%%%%%%%%%%%%%%%%%%%%%%%%%%%%%%%%%%%%%%%%%%%%%%%%%%%%%%%%%%%%%%%%%%%
% PREGLED NAREDBI PREDLOŠKA (sve iza \end{document} se ne kompajlira)
%
% Specijalni znakovi:  \$ \# \_ \{ \} \textasciitilde \textasciicircum \textbackslash
% Unicode znak po kodu: \char"0416
%
% Kratki listing (do ~10 linija, ne prelama se):
%   \begin{shortlisting} ... \end{shortlisting}
% Duži listing:
%   \begin{lstlisting} ... \end{lstlisting}
% Jezik listinga: \lstset{language=VHDL}  (ili Verilog, Python, C...)
%
% Nabrajanja: enumerate / itemize (predložak ih pravi bez velikih razmaka)
%
% Slika iz fajla: \langfigure{ime.png}{0.70}{Naslov slike}{fig:Ime}
%   (širina relativno u odnosu na papir; oznaka fig:...)
%   Referenca: Slika \ref{fig:Ime}
% TikZ crtež (profesor traži da se slike crtaju same, da ne izgledaju kopirane):
%   \begin{langfiguretikz}[]{0.3}{Naslov}{fig:Ime}
%     \draw (0,0) circle (1cm);
%     \node at (0,0) {Me};
%   \end{langfiguretikz}
%
% Jednostavna tabela (redovi odvojeni zarezom, kolone uspravnom crtom; prvi red je
% zaglavlje; | u polju = \vbr, zarez u polju = \com):
%   \langsimpletable{4}{Naslov tabele}{tab:Ime}
%   {
%   a    | b    | a \&\& b | a \vbr\vbr b ,
%   true | true | true     | true
%   }
% Složenija tabela: table + tabular s \rowcolor{lightgray} za zaglavlje
%   (primjer u originalnom predlošku).
%
% Gramatika:
%   \begin{grammar}
%     <A> $\rightarrow$ B | c | d
%   \end{grammar}
% Sintaksni dijagram: syntdiag okruženje, ponavljanje u rep, alternative u stack.
%
% Formula:
%   \begin{equation} \label{eqn:Ime}
%     y = f\left(\sum_{i=1}^{n} w_i x_i + b\right)
%   \end{equation}
%%%%%%%%%%%%%%%%%%%%%%%%%%%%%%%%%%%%%%%%%%%%%%%%%%%%%%%%%%%%%%%%%%%%%%%%%%%%%%%%
